\section{Deciding type equivalence}
\label{sec:type-equiv}

Type equivalence is decidable. 
\cref{fig:normalisation} contains the normalization algorithm. Its definition comprises two mutually
recursive functions $\Normal[+]{\cdot}$ and $\Normal[-]{\cdot}$. The 
superscript  is a polarity $+$ or $-$.
Function $\Normal[+]T$ is defined on all types, but 
function $\Normal[-]T$ is restricted to session types.
The negative polarity $-$ indicates that there is a pending $\TDual{\cdot}$
operation on the argument type, which is propagated by normalization.

Normalization leaves functional types unaffected. On session types,
the negative polarity is reified on type variables, it is propagated
to the session component of sending and receiving types, and each
traversal of a $\TDual\cdot$ inverts the polarity. Negative polarity also
inverts the direction of transmissions. However, the final direction is
determined by the polarity of the protocol.

\begin{example}
  Let's compute the positive normal form of
  $\TDual{(\TIn{(-\TInt)}{\alpha})}$ step by step.
  \begin{align*}
    & \Normal[+]{\TDual{(\TIn{-\TInt}{\alpha})}} \\
    ={}& \Normal[-]{{\TIn{-\TInt}{\alpha}}} \\
    ={}& \TFOut{\Normal[+]{-\TInt}}{\Normal[-]\alpha} \\
    ={}& \TFOut{\TFMinus{\Normal[+]{\TInt}}}{\TDual\alpha} \\
    ={}& \TFOut{\TFMinus{\TInt}}{\TDual\alpha} \\
    ={}& \TFOut{-{\TInt}}{\TDual\alpha} \\
    ={}& \TFIn{\TInt}{\TDual\alpha} \\
    ={}& \TIn{{\TInt}}{\TDual\alpha}
  \end{align*}
  In calculating $\Normal[-]{\typec{\TIn{-\TInt}{\alpha}}}$, the normal form of the
  protocol $\Normal[+]{-\TInt} = \typec{-\TInt}$ has negative polarity,
  which inverts the direction once more from output $\typec !$ to 
  input $\typec ?$.
\end{example}

To test whether type equivalence $\isConv TU\kind$ holds, we compute
the normal forms $\Normal[+]T$ and $\Normal[-]U$ and compare them
modulo $\alpha$ convertibility. 
We write $\typec T =_\alpha \typec U$ if types $\typec T$ and $\typec U$ are
$\alpha$-equivalent. The following propositions establish soundness and
completeness of this algorithm.

Every type is convertible to its normal form.
\begin{proposition}[Kind preservation]~\label{lemma:kind-preservation}
  \begin{enumerate}
  \item For all $\typeSynth T\kind$,  $\isConv{\Normal[+] T}T\kind$.
  \item For all $\typeSynth T\kinds$,  $\isConv{\Normal[-] T}{\TDual T}\kinds$.
  \end{enumerate}
\end{proposition}
\begin{proof}
  See Agda development \texttt{nf-sound+} and \texttt{nf-sound-}.
\end{proof}

\begin{corollary}[Soundness] Let $\typeSynth T\kind$ and $\typeSynth U\kind$.
  \begin{enumerate}
  \item If $\Normal[+] T =_\alpha \Normal[+] U$, 
    then $\isConv TU\kind$.
  \item If $\kind = \kinds$ and  $\Normal[-] T =_\alpha \Normal[-]
    U$, then $\isConv TU\kinds$.
  \end{enumerate}
\end{corollary}
\begin{proof}
  Item~1. Suppose that $\Normal[+] T =_\alpha \Normal[+] U$. By
  Lemma~\ref{lemma:kind-preservation}, we have  $\isConv{\Normal[+]
    T}T\kind$ and  $\isConv{\Normal[+] U}U\kind$. Conclude by symmetry
  and transitivity of conversion as ${=_\alpha} \subseteq {\equiv}$.

  The proof for item~2 is analogous.
\end{proof}

If two types are convertible, then their normal forms are equal.
\begin{proposition}[Completeness]   Let $\isConv TU\kind$. 
  \begin{enumerate}
  \item  $\Normal[+]T =_\alpha \Normal[+]U$.
  \item    If $\kind=\kinds$, then $\Normal[-]T =_\alpha \Normal[-]U$. 
  \end{enumerate}
\end{proposition}
\begin{proof}
  See Agda development \texttt{nf-complete} and \texttt{nf-complete-}.
\end{proof}

To obtain a complexity bound for the computation of a normal form, we
characterize the output of the function $\Normal T$, that is, the
syntax of normal forms. 
Normal forms for well formed types are $\typec{R^-}$ such that:
%
\begin{align*}
  \typec R^- \grmeq\;& \typec{R} \grmor \pneg{R}
  \\
  \typec{R} \grmeq\;&
  % functional types
  \TUnit \grmor \TFun[]{R}{R} \grmor \TPair{R}{R} \grmor
  \typec{\alpha} \grmor \TForall \alpha \kind {R}
  % session types
  \grmor \TEndW \grmor \TEndT \grmor \TIn{R}{R} \grmor \TOut{R}{R} \grmor \TDual{\alpha}
  % protocol and data types
  \grmor \TPApp X {\overline{R^-}}
\end{align*}
Essentially, the negative polarity can only occur at most once at the top level (of
a type of kind $\kindp$) and in the parameters of a protocol type. The
$\TDual\cdot$ operator only appears on type variables.
\begin{lemma}\label{lemma:syntax-of-type-normal-forms}
  If $\Normal[\pm]T = \typec{T'}$, then $\typec{T'}$ is described by $\typec{R^-}$.
\end{lemma}
\begin{proof}
  See Agda development \texttt{nf-normal}.
\end{proof}
\begin{proposition}
  For each $\typeSynth T\kind$, 
  the run time of $\Normal[\pm]T$ is in $O(|\typec{T}|)$.
\end{proposition}
\begin{proof}
  A run of $\Normal T$ visits each of the $|\typec{T}|$ nodes in
  the input type once. At each node, it either reconstructs the type
  from the normal forms of the sub-types in $O(1)$ time, or (in the transmission rules) it applies polarity
  correction to normal forms. By
  Lemma~\ref{lemma:syntax-of-type-normal-forms}, a normal form  
  starts with at most one negative polarity, so  correction
  takes at most two steps, that is $O(1)$. 
\end{proof}

%%% Local Variables:
%%% mode: latex
%%% TeX-master: "main"
%%% End:
